\documentclass[11pt,a4paper]{article}

\usepackage[utf8]{inputenc}
\usepackage[spanish,es-tabla]{babel}
\usepackage[margin=2cm]{geometry}
\usepackage{graphicx}
\usepackage{booktabs}
\usepackage{hyperref}
\usepackage{fancyhdr}

\pagestyle{fancy}
\fancyhf{}
\fancyhead[L]{\small Evidencias de Compilación y Ejecución - N-Body 2D}
\fancyhead[R]{\small Jonathan Zacarias (Ensamblado)}
\fancyfoot[C]{\thepage}

\title{
    \vspace{-1.5cm}
    \textbf{\Large Documento de Evidencia Experimental}\\[0.2cm]
    \textbf{\large Laboratorio PCAM + OpenMP: N-Body 2D}
}
\author{\textbf{Jonathan Zacarias, Hugo Barillas, June Herrera}\\ (Ensamblado por Jonathan Zacarias)}
\date{Octubre 2026}

\begin{document}

\maketitle

\section{Información del Hardware de Pruebas (Una Sola Máquina)}
\begin{itemize}
    \item \textbf{Procesador (CPU):} [Insertar modelo, ej. Intel Core i7 / AMD Ryzen]
    \item \textbf{Núcleos físicos / Lógicos:} [Ej. 8 núcleos / 16 hilos]
    \item \textbf{Memoria RAM:} [Ej. 16 GB DDR4]
    \item \textbf{Sistema Operativo:} [Ej. Ubuntu 22.04 LTS / Windows 11 WSL2]
    \item \textbf{Compilador:} \texttt{gcc (GCC) -O2 -fopenmp}
\end{itemize}

\section{Capturas de Compilación}
% \begin{figure}[h!]
% \centering
% \includegraphics[width=0.85\textwidth]{evidence/compilacion.png}
% \caption{Compilación limpia sin advertencias de las tres versiones con gcc -O2 -fopenmp}
% \end{figure}

\section{Capturas de Ejecución y Tiempos}

\subsection{Versión Secuencial (Baseline)}
% \begin{figure}[h!]
% \centering
% \includegraphics[width=0.85\textwidth]{evidence/ejecucion_secuencial.png}
% \caption{Ejecución de nbody_secuencial con N=5000, 10 pasos, seed=42}
% \end{figure}

\subsection{Versión Paralela (OpenMP)}
% \begin{figure}[h!]
% \centering
% \includegraphics[width=0.85\textwidth]{evidence/ejecucion_paralelo.png}
% \caption{Ejecución de nbody_paralelo con 1, 2, 4 y 8 hilos}
% \end{figure}

\subsection{Versión Optimizada}
% \begin{figure}[h!]
% \centering
% \includegraphics[width=0.85\textwidth]{evidence/ejecucion_optimizada.png}
% \caption{Ejecución de nbody_optimized validando el Checksum}
% \end{figure}

\section{Tabla Resumen de Experimentos}
% \begin{table}[h!]
\centering
\caption{Tiempos promedio (5 repeticiones), speedup y eficiencia. $N=5000$, 10 pasos, $T_1 = 1.788$ s}
\label{tab:experimentos}
\small
\begin{tabular}{@{}lcccrrrrr@{}}
\toprule
\textbf{Versión} & \textbf{$p$} & \textbf{Schedule} & \textbf{Chunk} & \textbf{$\bar{T}$ (s)} & \textbf{$\sigma$ (s)} & \textbf{$S(p)$} & \textbf{$E(p)$} & \textbf{$e_{KF}$} \\ \midrule
nbody\_secuencial & 1 & sequential & default & 1.788 & 0.245 & 1.00 & 1.00 & -- \\
nbody\_paralelo & 1 & static & default & 1.781 & 0.168 & 1.00 & 1.00 & -- \\
nbody\_paralelo & 2 & static & default & 1.082 & 0.087 & 1.65 & 0.83 & 0.211 \\
nbody\_paralelo & 4 & static & default & 0.748 & 0.054 & 2.39 & 0.60 & 0.225 \\
nbody\_paralelo & 8 & static & default & 0.547 & 0.025 & 3.27 & 0.41 & 0.207 \\
nbody\_paralelo & 8 & static & 8 & 0.542 & 0.031 & 3.30 & 0.41 & 0.204 \\
nbody\_paralelo & 8 & static & 64 & 0.570 & 0.048 & 3.14 & 0.39 & 0.222 \\
nbody\_paralelo & 8 & dynamic & 8 & 0.474 & 0.028 & 3.77 & 0.47 & 0.160 \\
nbody\_paralelo & 8 & dynamic & 64 & 0.519 & 0.025 & 3.45 & 0.43 & 0.189 \\
nbody\_paralelo & 8 & guided & default & 0.537 & 0.049 & 3.33 & 0.42 & 0.200 \\
\bottomrule
\end{tabular}
\end{table}
 o pegar tabla de results.csv

\end{document}
